\section{Threat Model and Classification}\label{sec:taxonomy}
\textbf{A result is meaningful only for the target and observations actually supplied to its attacker.}
The ticket example has client--Agent, Agent--model, Agent--tool and coordination links (Figure~\ref{fig:chain}). A transit observer sees connection destinations, directions, lengths and times. A contacted service reads the request addressed to it; an active auditor can also choose challenges. Table~\ref{tab:deployment-access} maps these rights to deployment configurations. The same traffic may encode requested intent, a logged action or successful execution. Authorization additionally requires the principal and policy context: identifying a refund does not decide whether it was approved.

\begin{table*}[t]\centering\small
\caption{Deployment and access map. The configurations are analyst-defined; cited settings test the listed view, not every component of the deployment. A gateway's plaintext access requires TLS termination.}\label{tab:deployment-access}
\setlength{\tabcolsep}{4pt}
\begin{tabularx}{\textwidth}{@{}P{.18\textwidth}P{.25\textwidth}P{.28\textwidth}Y@{}}\toprule
Deployment & Observer and permitted view & Existing evidence for that view & Protection consequence \\\midrule
Enterprise egress & A transit monitor sees routed model/tool flows, sizes, times and any exposed destination association. & WRA: supplied research-domain traces; ANT and cloud benchmark: model/combined traffic~\cite{jeong2026wra,li2026ant,pouryousef2026traffic}. & Specify which services cross the monitored egress. A union of flows is available only when the monitor can observe them. \\\addlinespace[3pt]
Shared gateway or relay & Transit sees multiplexed gateway output. A terminating gateway can additionally read addressed requests. & Live OHTTP: passive-link inference under stated relay trust~\cite{dangol2026workflow}; PrivacyTrace: semantic recipient/gateway access~\cite{duan2026privacytrace}. & Separate relay association protection from content trust. Connection sharing changes the observation unit. \\\addlinespace[3pt]
Local model, remote tools & An external observer sees remote tool traffic; the local model call has no corresponding external network leg. & WRA's destination-recovery mechanism is relevant to tool egress~\cite{jeong2026wra}; whole-deployment protection is untested in the mapped comparison. & Evaluate tool and downstream calls. Do not import a remote-model attack score into an absent network leg. \\\addlinespace[3pt]
Cloud model, remote tools & An authorized vantage may see either or both encrypted legs. A contacted tool sees its own arguments. & ANT/cloud benchmark test model traffic~\cite{li2026ant,pouryousef2026traffic}; PrivacyTrace tests recipient-visible representations~\cite{duan2026privacytrace}. & Check permitted model and tool views after the intervention; recipients and passive monitors need different protections. \\\bottomrule
\end{tabularx}
\end{table*}


\subsection{What is observed and what is predicted}
A captured packet, TLS record, HTTP stream, tool call and complete task are different units. A gateway can regroup streamed output; HTTP/2 and QUIC can multiplex streams; MCP may use local stdio or remote HTTP~\cite{mcp2025transports,rfc9113,rfc9114,rfc9000,rfc8446}. A boundary supplied by the evaluator is prior information, not a classifier output. Six fields specify each inference setting: interface and representation, observer access, supplied prior information, prediction target, evaluation population and base rate, and protected-training access. Host traces can supply experimental truth without becoming transit-visible features~\cite{reynolds2006wap5,tak2009vpath}.

\subsection{When cover must execute}\label{sec:necessity}
Transport cover ends before a business tool, at an endpoint the defender controls~\cite{apthorpe2019stp}. That endpoint must sit after the observer: otherwise the observer can see it discard the cover. A defender can sometimes move the endpoint, for example to a relay outside a monitored egress, but the segment from the last defender-controlled hop to the service remains. If the adversary can observe that segment, or is the service itself, every request it sees must be real. Cover then hides a real call only by arriving as a real call (Table~\ref{tab:cover-necessity}). Executed cover therefore follows from the deployment; it is not a design preference. It crosses an application boundary: it may consume quota, access other objects, update state or enter the planner's context.

\begin{table}[t]\centering\footnotesize
\caption{When cover must execute. Transport cover needs a defender-controlled terminator located after the observer and before the business service. Without one, cover can hide the real call only by arriving as a real call.}\label{tab:cover-necessity}
\setlength{\tabcolsep}{3pt}
\begin{tabularx}{\columnwidth}{@{}P{.27\columnwidth}P{.2\columnwidth}Y@{}}\toprule
Observer $\to$ service & Terminator after observer? & Cover that can work \\\midrule
Egress $\to$ own service & Yes (own tunnel) & Transport cover suffices~\cite{apthorpe2019stp} \\\addlinespace[2pt]
Egress $\to$ third-party SaaS or remote MCP & No & Executed cover; isolation only if the provider offers it \\\addlinespace[2pt]
Recipient (tool provider) & No & Executed cover whose arguments must also be plausible to the recipient \\\addlinespace[2pt]
TLS-terminating gateway & After gateway only & Executed cover, or trust the gateway with content \\\bottomrule
\end{tabularx}\end{table}


\subsection{A taxonomy of cover}\label{sec:cover-taxonomy}
We classify a cover policy along four axes (Table~\ref{tab:cover-taxonomy}). \emph{Termination and state} determine what business work cover performs. \emph{Feedback visibility} determines whether cover responses enter later model requests. The \emph{dispatcher}, the process that decides when cover is sent, determines whether cover inherits the timing of honest calls. Honest calls follow a model decision; a controller that inserts cover without waiting for the model creates a timing signature that has no counterpart in the first three axes. The live study varies the dispatcher while holding termination (an isolated clone) and feedback (hidden) fixed.

\begin{table*}[t]\centering\small
\caption{Cover taxonomy. Each policy is classified by where it terminates and what state it touches, whether its responses reach the real planner, and which process dispatches it. The last column names the evaluation in which the policy is measured; S1 is the earlier state control and Section~\ref{sec:cover-study} the live timing study.}\label{tab:cover-taxonomy}
\setlength{\tabcolsep}{4pt}
\begin{tabularx}{\textwidth}{@{}P{.17\textwidth}P{.26\textwidth}P{.13\textwidth}P{.2\textwidth}Y@{}}\toprule
Policy & Terminates at / state touched & Feedback to planner & Dispatcher & Measured in \\\midrule
Transport cover & Tunnel endpoint before the service; no business state & No & Sender schedule & STP~\cite{apthorpe2019stp}; WF padding~\cite{juarez2016wtfpad,gong2020front} \\
Proxy discard & Defender proxy computes on a disposable copy & No & Controller & S1 \\
Read cover & Real service; reads only & No & Controller & S1 \\
Isolated execution & Real service on a persistent clone & No / yes & Controller & S1 (hidden and visible); all live cover arms \\
Live write & Real service on production state & No & Controller & S1 (stress control) \\
Research decoys & Real third-party Web destinations & Not located & Decoy generator & WRA~\cite{jeong2026wra} \\\midrule
Immediate & Isolated clone & No & Controller, directly after an honest response & Section~\ref{sec:cover-study} (calibration) \\
History-matched & Isolated clone & No & Controller, training-gap and duration samples; sequential or concurrent & Section~\ref{sec:cover-study} \\
Planner shadow & Isolated clone; separate planner context & No & Second planner on an independent ticket & Section~\ref{sec:cover-study} \\\bottomrule
\end{tabularx}\end{table*}


The live study models the second row of Table~\ref{tab:cover-necessity} when the observed hop is the last one before the service: seven MCP services stand for third-party tools, no defender relay intervenes, and the observer is a transit monitor on that path. Information forwarded to another service creates another permitted observation; it does not imply that every service is visible to the same attacker. The recipient view of executed cover, where the attacker also reads decoy arguments, is not evaluated here and remains open (Section~\ref{sec:agenda}).
